% !TEX root = ../main.tex
% Student declaration (UNISA wording). The candidate must review the AI-use
% statement, then sign and date before submission.

\clearpage
\phantomsection
\chapter*{선언}
\label{ch:declaration}

\noindent 이름: 권태홍 (Taehong Kwon)\\
학번: 28576810\\
학위: 컴퓨터과학 박사

\medskip
\noindent ERC-20 허락 화살표를 페이지랭크에 넣어 온체인 지갑 평판 점수를 더 좋게 만들기

\medskip
저는 위의 논문이 제가 직접 한 일이라고 선언해요. 그리고 제가 쓰거나 그대로 따온 자료는 모두 빠짐없는 참고문헌으로 표시하고 출처를 밝혔다고 선언해요.

또 저는 이 논문을 독창성 검사 프로그램(다른 글을 베끼지 않았는지 살펴보는 프로그램)에 넣어 검사받았다고 선언해요. 그리고 그 결과가 독창성에 대해 받아들이는 기준 안에 든다고 선언해요.

또 저는 이 연구를 전에 심사받으려고 다음 두 곳에 낸 적이 없다고 선언해요. 연구 전체도, 그 일부도 마찬가지예요. 다른 자격을 받으려고 Unisa(남아프리카대학교)에 낸 적이 없고, 다른 고등 교육 기관(대학 같은 곳)에 낸 적도 없어요.

\paragraph{생성형 AI 도구를 쓴 일.} 생성형 AI 도구(글이나 코드를 만들어 주는 인공지능 도구)인 Claude(Anthropic)를 다음 일을 돕는 데 썼어요. 문장 다듬기, 글의 짜임을 고치고 내용이 앞뒤로 맞는지 서로 맞춰 보기, 분석 스크립트(분석에 쓰는 짧은 프로그램)의 일부를 쓰고 시험해 보기, 그리고 참고할 글을 찾고 확인하기예요. 연구 질문, 연구 설계, 자료, 미리 정한 판정 규칙, 분석, 결론은 모두 제 것이에요. 저는 논문에 적은 숫자를 하나도 빠짐없이 공개 저장소에 있는 요약 파일(기계가 읽을 수 있는 파일)과 맞춰 보았어요. 그리고 이 논문의 내용은 모두 제가 책임져요.


\vspace{2em}
\noindent\rule{6cm}{0.4pt}\hfill\rule{4cm}{0.4pt}\\
서명 \hfill 날짜
